
\documentclass[a4paper,11pt]{article}
\usepackage{amsmath}
\begin{document}
\section{Mathematische Formeln}

Das ist ein guter Versuch um mathematische formeln zu verwenden. 
Hier versuchen wir es einfach mal Inline mit einer einfachen Formel wie dieser 
oder anderen: \begin{math}  a^{2} + b^{2} = c^{2}\end{math} (hier ist es der Satz des Pythagoras). \\
\\
\\

Man sollte aber auch wissen, dass Manche Sachen besser als eigenen Zeile stehen sollten dann nutzt man 
\begin{displaymath}
 a^{2} + b^{2} = c^{2}
\end{displaymath}
Damit man es besser lesen kann. 
Die Fakultät einer Zahl n ist definiert als \begin{math}n! = 1 \times 2 \times \ldots \times n.\end{math} Mathematiker definieren
das präziser als
\begin{displaymath}
n! =  \prod_{i = 1}^{x}{i}
\end{displaymath}
Aber das ist sowieso unvollständig, da ja auch $n = 0$ gelten kann. Richtig vollständig
muß man also eine Fallunterscheidung treffen und kann bei der Gelegenheit das ganze
auch rekursiv definieren.
\begin{displaymath}
\text{Fak}(n) = 
\begin{cases}
1  &  \text{falls } n = 0  \\
n * \text{Fak}(n - 1)  &  \text{sonst} 
\end{cases}
\end{displaymath}
Man beachte das in dieser Formel ver¨anderte Multiplikationssymbol. Außerdem benutzen
Mathematiker auch gerne griechische Buchstaben, wie $\alpha, \beta, \gamma$. Wenn dann die Buchstaben
knapp werden, dann wird Hoch- oder Tiefgestellt, so wie in $x_{1}, x_{2}, \ldots , x_{n} \text{ oder } y^{2}, y^{3}.$
\end{document}